\documentclass[10pt,a4paper]{amsart}
\usepackage[margin=19mm]{geometry}
\usepackage{amsmath,amssymb,mathtools}
\usepackage[scr=rsfs,cal=euler]{mathalfa}
\usepackage{xfrac}
\usepackage[unicode,hidelinks,bookmarks=false]{hyperref}
\newcommand{\ie}{i.e.\ }
\newcommand{\cf}{cf.\ }
\newcommand{\resp}{respectively\ }
\newcommand{\Subsection}{\subsection}
\providecommand{\nobreakdash}{\nobreak\hbox{-}}
\setlength{\emergencystretch}{2em}
\multlinegap0pt
\usepackage{bm}
\usepackage{bbm}
\usepackage{dsfont}
\usepackage{xspace}
\usepackage{cancel}
\usepackage{mathtools}
\usepackage{amsthm}
\makeatletter
\@ifundefined{thm}{\newtheorem{thm}{Theorem}}{}
\makeatother
\title[The classical Yangian symmetry of Auxiliary Field Sigma Mode]{The classical Yangian symmetry of Auxiliary Field Sigma Models}
\author{Daniele Bielli, Christian Ferko, Michele Galli, Gabriele Tartaglino-Mazzucchelli}
\date{Source: arXiv:2605.18213v1 (CC BY 4.0)}
\begin{document}
\maketitle

\noindent\emph{Source and licence.} The classical Yangian symmetry of Auxiliary Field Sigma Models. Version arXiv:2605.18213v1 (CC BY 4.0), distributed under the Creative Commons Attribution 4.0 International licence, as recorded in the arXiv licence metadata for this exact version and shown on the abstract page https://arxiv.org/abs/2605.18213v1. Full rights notice: licenses/arxiv-2605.18213-CC-BY-4.0.txt.\par\smallskip

\noindent\textit{Editorial scope.} Counted unit: the Thm whose source label is \texttt{gBIZZ-theorem} in the article (source0.tex lines 417--453, source label \texttt{gBIZZ-theorem}), delivered together with the article's own complete original proof of it (source0.tex lines 455--584, one \texttt{proof} environment of 130 lines). No mathematics is written, repaired or re-derived: statement and proof are the article's own text line for line. Required context retained uncounted: j-properties-definition (lines 310--317); original-BIZZ-tower (lines 319--328); BIZZ-potentials (lines 330--336); form-commutator-star-identities (lines 2728--2736); useful-Jacobi-1 (lines 2738--2743); useful-Jacobi-2 (lines 2745--2749). Every definition, display and label of those lines is retained unchanged. Omitted from this package and not redistributed: 1--309, 318--318, 329--329, 337--416, 454--454, 585--2727, 2737--2737, 2744--2744, 2750--3930. Nothing inside a retained excerpt is omitted or altered: every retained body is delivered exactly as the article's lines are written, so no omission receipt of any kind accompanies this package. Citations: the retained excerpts carry no citation command at all, so no bibliography entry of the article is transcribed and no reference is redistributed. Reference adaptation: none was needed and none was made; every reference inside the delivered unit resolves inside it, so no unresolved reference or citation is produced. Printed slips kept literal: no printed word, symbol, hypothesis or number of the counted statement or of its proof was altered. Layout and class: the article's own class is \texttt{article} with packages including xcolor, verbatim, tensor, amsmath, amssymb, calc, amsthm, bbm, epsfig, psfrag, mathtools, latexsym, bm, amsfonts; the wrapper is the lane's pinned \texttt{amsart} editorial wrapper at 10pt on A4, which restates the theorem-like environments the retained text needs and adds the article's own macro definitions that the retained text needs (none), with extra wrapper packages (bm, bbm, dsfont, xspace, cancel, mathtools). The delivered numbering is the unit's own. Assets: no retained span contains a figure environment, a graphics-inclusion command, a PDF-page inclusion, a table or a TikZ picture; no image file is copied into this package, the package declares no asset dependency and no image file is redistributed with it. Licence: version arXiv:2605.18213v1 of this article is distributed under the Creative Commons Attribution 4.0 International licence, as recorded in the arXiv licence metadata for this exact version and shown on the abstract page https://arxiv.org/abs/2605.18213v1. A full rights notice is delivered at licenses/arxiv-2605.18213-CC-BY-4.0.txt. Characters: every retained excerpt is pure ASCII, so the delivered engine, XeLaTeX with its default Latin Modern text font, reports no missing character.
\begin{equation}\label{j-properties-definition}
\begin{aligned}
\text{Conservation:}& \qquad
0=\mathrm{d}(\star L) \, ,
\\
\text{Flatness:}& \qquad 0=\mathrm{d}L+\tfrac{1}{2}[L,L] \,\, .
\end{aligned}
\end{equation}

\begin{equation}\label{original-BIZZ-tower}
\begin{aligned}
J^{(0)}&:=\beta \, L
\quad \quad \quad \quad \,\,\, \forall \, \beta \in \mathbb{R} \, ,
\\
J^{(1)}&:=\star L+\tfrac{1}{2\beta}[L,\chi^{(0)}] \, ,
\\
J^{(n)}&:=\nabla_{L}\chi^{(n-1)} \qquad \forall \, n \geq 2 \qquad \text{with} \qquad \nabla_{L}:=\mathrm{d}+[L,-] \,\, ,
\end{aligned}
\end{equation}

\begin{equation}\label{BIZZ-potentials}
\mathrm{d}(\star J^{(i)})=0 
\qquad \Longleftrightarrow \qquad
\star J^{(i)}=\mathrm{d}\chi^{(i)} \,\, , 
\qquad \qquad \forall \, i\geq 0
\,\, .
\end{equation}

\begin{equation}\label{form-commutator-star-identities}
\begin{aligned}
[\alpha,\star\beta]&=-[\star\alpha,\beta]=-[\beta,\star\alpha] 
\qquad , \qquad 
[\alpha,\star\alpha]=-[\star\alpha,\alpha]=-[\alpha,\star \alpha]=0 \,\, ,
\\
&\qquad \text{and} \quad\qquad [\star \alpha,\star \beta]=-[\star^2\alpha,\beta]=-[\alpha,\beta] \,\, .
\end{aligned}
\end{equation}

\begin{equation}\label{useful-Jacobi-1}
[\alpha,[\alpha,f]]=-\tfrac{1}{2}[f,[\alpha,\alpha]] 
\quad \text{for} \quad \beta=\alpha\in \Omega^{1}(\Sigma,\mathfrak{g})
\quad \text{and} \quad 
f\in \Omega^{0}(\Sigma,\mathfrak{g}) \,\, ,
\end{equation}

\begin{equation}\label{useful-Jacobi-2}
[\alpha,[\star \beta,f]]=-[\star \beta,[\alpha,f]] \,\, ,
\qquad \text{for} \qquad \alpha,\beta\in \Omega^{1}(\Sigma,\mathfrak{g}) \qquad \text{and} \qquad
[\alpha,\star\beta]\equiv0 \,\, .
\end{equation}

\begin{thm}\label{gBIZZ-theorem}
Consider two Lie-algebra valued 1-forms $A$ and $B$ satisfying
\begin{equation}\label{A,B-definition}
\begin{aligned}
\text{Conservation:}& \qquad
0=\mathrm{d}(\star B)  \, ,
\\
\text{Flatness:}& \qquad 0=\mathrm{d}A+\tfrac{1}{2}[A,A] \, ,
\end{aligned}
\end{equation}
and commutation relations
\begin{equation}\label{A,B-commutators}
[A,A]=[B,B]
\qquad \text{and} \qquad 
[A,\star B]=[B,\star A]=0 \,\, .
\end{equation}
Then the following currents generalise the BIZZ tower \eqref{original-BIZZ-tower} and are conserved:
\begin{equation}\label{generalised-BIZZ-tower}
\begin{aligned}
J^{(0)}&:=\beta \, B \qquad \qquad \qquad \qquad \qquad \qquad \quad \quad \,\, \forall \, \beta \in \mathbb{R}\,\, ,
\\
J^{(1)}&:= \star A + \tfrac{1}{2\beta}[B,\chi^{(0)}] \, ,
\\
J^{(2)}&:=B+\tfrac{1}{2\beta}[\star A,\chi^{(0)}]+[B,\chi^{(1)}] \, ,
\\
J^{(n)}&:=J^{(n-2)}+[\star A,\chi^{(n-2)}]+[B,\chi^{(n-1)}] \qquad \forall \, n\geq 3 \,\, ,
\end{aligned}
\end{equation}
with $\chi^{(i)}$ local potentials defined as in \eqref{BIZZ-potentials}
\begin{equation}\label{gBIZZ-potentials}
\mathrm{d}(\star J^{(i)})=0 
\qquad \Longleftrightarrow \qquad
\star J^{(i)}=\mathrm{d}\chi^{(i)} \,\, , 
\qquad \qquad \forall \, i\geq 0
\,\, .
\end{equation}
\end{thm}

\begin{proof} It is straightforward to recognise that for $A=B=L$ the commutators \eqref{A,B-commutators} are identically satisfied and the conditions \eqref{A,B-definition} become \eqref{j-properties-definition}, such that the tower \eqref{generalised-BIZZ-tower} reduces to \eqref{original-BIZZ-tower}. The proof of conservation can be divided into two sectors.

\textbf{Conservation of $J^{(0)},J^{(1)},J^{(2)}$.} \qquad This can be verified by direct computation: for $J^{(0)}$ it follows trivially from \eqref{A,B-definition}, while for $J^{(1)}$ one finds 
\begin{equation}
\begin{aligned}
\mathrm{d}(\star J^{(1)})&=\mathrm{d}(A+\tfrac{1}{2\beta}[\star B,\chi^{(0)}])
\\
&=\mathrm{d}A+\tfrac{1}{2\beta}[\mathrm{d}(\star B),\chi^{(0)}]-\tfrac{1}{2\beta}[\star B,\mathrm{d}\chi^{(0)}]
\\
&=-\tfrac{1}{2}[A,A]-\tfrac{1}{2\beta}[\star B,\star J^{(0)}]
\\
&=-\tfrac{1}{2}[A,A]+\tfrac{1}{2\beta}[ B, J^{(0)}]
\\
&=-\tfrac{1}{2}[A,A]+\tfrac{1}{2}[ B, B]
\\
&=0 \,\, ,
\end{aligned}
\end{equation}
where in going from the third to the fourth line we picked up a sign on the second term using \eqref{form-commutator-star-identities} and in the last step the two terms cancel due to \eqref{A,B-commutators}. Similarly, for $J^{(2)}$ 
\begin{equation}
\begin{aligned}
\mathrm{d}(\star J^{(2)})&=\mathrm{d}(\star B+\tfrac{1}{2\beta}[A,\chi^{(0)}]+[\star B,\chi^{(1)}])
\\
&=\tfrac{1}{2\beta}[\mathrm{d}A,\chi^{(0)}]-\tfrac{1}{2\beta}[A,\mathrm{d}\chi^{(0)}]-[\star B,\mathrm{d}\chi^{(1)}]
\\
&=\tfrac{1}{2\beta}[\mathrm{d}A,\chi^{(0)}]-\tfrac{1}{2\beta}[A,\star J^{(0)}]-[\star B,\star J^{(1)}]
\\
&=\tfrac{1}{2\beta}[\mathrm{d}A,\chi^{(0)}]-\tfrac{1}{2\beta}[A,\star J^{(0)}]+[ B, J^{(1)}]
\\
&=-\tfrac{1}{4\beta}\bigl[[A,A],\chi^{(0)}\bigr]-\tfrac{1}{2}[A,\star B]+[B,\star A]+\tfrac{1}{2\beta}\bigl[B,[B,\chi^{(0)}]\bigr]
\\
&=\tfrac{1}{4\beta}\bigl[\chi^{(0)},[A,A]\bigr]+\tfrac{1}{2\beta}\bigl[B,[B,\chi^{(0)}]\bigr]
\\
&=0 \,\, ,
\end{aligned}
\end{equation}
where in the first line we exploited conservation of $B$, in the second the definition of the potentials \eqref{BIZZ-potentials}, in the third the identities \eqref{form-commutator-star-identities}, in the fourth the commutators \eqref{A,B-commutators}, and in the last again \eqref{A,B-commutators} and the Jacobi identity \eqref{useful-Jacobi-1}.


\textbf{Conservation of $J^{(n)}$ for $n\geq3$.} \qquad To show conservation of the remaining currents in the tower we proceed in three steps. First, notice that conservation for $n\geq3$ reads
\begin{equation}
\begin{aligned}
\mathrm{d}(\star J^{(n)})= [\star A,J^{(n-2)}]+[B,J^{(n-1)}]+\tfrac{1}{2}\bigl[\chi^{(n-2)},[A,A]\bigr] =: X^{(n-1)} \,\, .
\end{aligned}
\end{equation}
By direct computation one indeed finds that
\begin{equation}
\begin{aligned}
\mathrm{d}(\star J^{(n)})&=\mathrm{d}(\star J^{(n-2)}+[A,\chi^{(n-2)}]+[\star B,\chi^{(n-1)}])
\\
&=[\mathrm{d}A,\chi^{(n-2)}]-[A,\mathrm{d}\chi^{(n-2)}]-[\star B,\mathrm{d}\chi^{(n-1)}]
\\
&=-\tfrac{1}{2}\bigl[ [A,A],\chi^{(n-2)}\bigr]-[A,\star J^{(n-2)}]-[\star B,\star J^{(n-1)}]
\\
&=\tfrac{1}{2}\bigl[\chi^{(n-2)},[A,A]\bigr]+[\star A, J^{(n-2)}]+[B, J^{(n-1)}]
\\
&=: X^{(n-1)} \,\, ,
\end{aligned}
\end{equation}
where we rearranged using \eqref{BIZZ-potentials}, flatness \eqref{A,B-definition} of $A$, and commutator identities \eqref{form-commutator-star-identities}. The second step consists of noting that, thanks to the recursive definition of the currents for $n \geq 3$, the conservation conditions are recursively related as
\begin{equation}
X^{(n+1)}=X^{(n-1)} \qquad \forall \, n\geq 3 \,\, .
\end{equation}
This can once more be checked explicitly:
\begin{equation}
\begin{aligned}
X^{(n+1)}=&+[\star A,J^{(n)}]+[B,J^{(n+1)}]+\tfrac{1}{2}\bigl[\chi^{(n)},[A,A]\bigr] 
\\
=&+[\star A,J^{(n-2)}]+\bigl[\star A,[\star A,\chi^{(n-2)}] \bigr]+\bigl[\star A,[B,\chi^{(n-1)}]  \bigr]
\\
&+[B,J^{(n-1)}]+\bigl[B,[\star A,\chi^{(n-1)}] \bigr]+\bigl[B,[B,\chi^{(n)}] \bigr]+\tfrac{1}{2}\bigl[\chi^{(n)},[A,A]\bigr]
\\
=&+\bigl[\star A,[B,\chi^{(n-1)}]  \bigr]+\bigl[B,[\star A,\chi^{(n-1)}] \bigr]
\\
&+\bigl[B,[B,\chi^{(n)}] \bigr]+\tfrac{1}{2}\bigl[\chi^{(n)},[A,A]\bigr]+X^{(n-1)}
\\
=&X^{(n-1)} \,\, .
\end{aligned}
\end{equation}
In the first line we substituted the recursive definition of $J^{(n)}$ and $J^{(n+1)}$ and in the second we collected the terms defining $X^{(n-1)}$, after noting that 
\begin{equation}
\bigl[\star A,[\star A,\chi^{(n-2)}] \bigr]=-\bigl[ A,[ A,\chi^{(n-2)}] \bigr] =+\tfrac{1}{2}\bigl[ \chi^{(n-2)},[A,A]\bigr] \,\, ,
\end{equation}
as a result of \eqref{form-commutator-star-identities} and the Jacobi identity \eqref{useful-Jacobi-1}. In going to the last line we exploited that the four extra terms cancel pairwise as a result of the Jacobi identities \eqref{useful-Jacobi-1}, \eqref{useful-Jacobi-2}, and the commutator identity \eqref{A,B-commutators}. At this point we have shown that as a result of the recursive definition \eqref{generalised-BIZZ-tower} for currents $J^{(n)}$ with $n\geq 3$, their conservation conditions also exhibit a recursive structure of the form
\begin{equation}
\mathrm{d}(\star J^{(n)})=X^{(n-1)}=X^{(n-3)}=X^{(n-5)} = \ldots \, ,
\end{equation}
and to show conservation of the whole tower it is sufficient to check the vanishing of the first non-trivial even and odd cases, namely $X^{(2)}=0$ and $X^{(3)}=0$, which constitutes the third and final step of the proof. For $X^{(2)}$ one finds that
%
\begin{equation}
\begin{aligned}
X^{(2)}=&+[\star A,J^{(1)}]+[ B,J^{(2)}]+\tfrac{1}{2}\bigl[\chi^{(1)},[A,A]\bigr]
\\
=&+[\star A,\star A]+\tfrac{1}{2\beta}\bigl[\star A,[B,\chi^{(0)}]\bigr]
\\
&+[B,B]+\tfrac{1}{2\beta}\bigl[B,[\star A,\chi^{(0)}]\bigr]+\bigl[B,[B,\chi^{(1)}]\bigr]+\tfrac{1}{2}\bigl[\chi^{(1)},[A,A]\bigr]
\\
=&-[A,A]+[B,B]+
\\
&+\tfrac{1}{2\beta}\bigl[\star A,[B,\chi^{(0)}]\bigr]+\tfrac{1}{2\beta}\bigl[B,[\star A,\chi^{(0)}]\bigr]
\\
&-\tfrac{1}{2}\bigl[\chi^{(1)},[B,B]\bigr]+\tfrac{1}{2}\bigl[\chi^{(1)},[A,A]\bigr]
\\
=&0 \,\, ,
\end{aligned}
\end{equation}
where in the first line we substituted the definitions \eqref{generalised-BIZZ-tower} of $J^{(1)}$ and $J^{(2)}$, in the second we exploited the identities \eqref{form-commutator-star-identities}, and in the third we noticed that terms cancel pairwise after using \eqref{A,B-commutators} and the Jacobi identity \eqref{useful-Jacobi-2}. Similarly, for $X^{(3)}$ one finds
\begin{equation}
\begin{aligned}
X^{(3)}=&+[\star A,J^{(2)}]+[B,J^{(3)}]+\tfrac{1}{2}\bigl[\chi^{(2)},[A,A]\bigr]
\\
=&+[\star A,B]+\tfrac{1}{2\beta}\bigl[\star A,[\star A,\chi^{(0)}]\bigr]+\bigl[\star A,[B,\chi^{(1)}]\bigr]+\tfrac{1}{2}\bigl[\chi^{(2)},[A,A]\bigr]
\\
&+[B,J^{(1)}]+\bigl[B,[\star A,\chi^{(1)}]\bigr]+\bigl[B,[B,\chi^{(2)}]\bigr]
\\
=&-\tfrac{1}{2\beta}\bigl[A,[A,\chi^{(0)}]\bigr]+\bigl[\star A,[B,\chi^{(1)}]\bigr]+\tfrac{1}{2}\bigl[\chi^{(2)},[A,A]\bigr]
\\
&+[B,\star A]+\tfrac{1}{2\beta}\bigl[B,[B,\chi^{(0)}]\bigr]+\bigl[B,[\star A,\chi^{(1)}]\bigr]+\bigl[B,[B,\chi^{(2)}]\bigr]
\\
=&+\tfrac{1}{4\beta}\bigl[\chi^{(0)},[A,A]\bigr]-\tfrac{1}{4\beta}\bigl[\chi^{(0)},[B,B]\bigr]+
\\
&+\bigl[\star A,[B,\chi^{(1)}]\bigr]+\bigl[B,[\star A,\chi^{(1)}]\bigr]
\\
&+\tfrac{1}{2}\bigl[\chi^{(2)},[A,A]\bigr]-\tfrac{1}{2}\bigl[\chi^{(2)},[B,B]\bigr]
\\
=&0 \,\, ,
\end{aligned}
\end{equation}
where in the first line we substituted the definitions of $J^{(2)}$ and $J^{(3)}$ in \eqref{generalised-BIZZ-tower}, and in the second we substituted the expression for $J^{(1)}$ in \eqref{generalised-BIZZ-tower} and simplified a few terms using \eqref{A,B-commutators}, \eqref{form-commutator-star-identities}. In the third line we rearranged terms using the Jacobi identity \eqref{useful-Jacobi-1} and in the fourth one we just exploited cancellations resulting from \eqref{A,B-commutators} and \eqref{useful-Jacobi-2}.
\end{proof}

\end{document}
